%%
%% test_comserifa.tex -- guarda a opcao de classe `comserifa'.
%%
%% A 2.2(b) do manual fixa a cor, o corpo 12 e o corpo menor uniforme das
%% citacoes longas, notas, paginacao e legendas, e NAO prescreve familia de
%% fonte; a Norma da COPPE tambem nao. Desde a v4.1 a classe compoe SEM SERIFA
%% por padrao, como o proprio manual do SiBI, que carrega a regra e esta
%% composto em Arial. `comserifa' existe para o autor que prefere a serifada.
%%
%% O teste pede as duas coisas ao mesmo tempo, `comserifa' e `pdfa', porque a
%% pergunta que interessa nao e se compila: e se o documento com serifa
%% continua sendo PDF/A. Sem fonte vetorial nao ha PDF/A, e foi exatamente
%% assim que a classe inteira saia antes de carregar o lmodern: em Type 3.
%% O veraPDF deste alvo esta no escopo `pdfa' do build-check.
%%
%% Cobre tambem os cinco niveis de secao, a citacao longa, a legenda e a nota
%% de rodape -- os quatro itens que a 2.2(b) manda sair em corpo menor -- para
%% que a troca de familia nao os deixe para tras.
%%
%% NAO e gerado pelo coppe.dtx e NAO e distribuido: a suite prova que a
%% classe funciona, nao faz parte dela. Rode por tools\prova.ps1, ou por
%% tests\run-tests.ps1 para so esta pasta.
%%

\documentclass[dsc,comserifa,pdfa]{coppe}

\usepackage[brazilian]{babel}

\title{Título de Teste com Serifa}
\foreigntitle{Serif Test Title}
\author{Nome}{de Teste}

\advisor{Primeiro}{Orientador}{D.Sc.}{UFRJ}
\examiner{Primeiro Examinador}{D.Sc.}{UFRJ}
\examiner{Segundo Examinador}{Ph.D.}{UFRJ}

\department{PESC}
\date{09}{2026}
\dataaprovacao{15 de setembro de 2026}

\keyword{Com serifa}
\keyword{Tipografia}

\areaconcentracao{Engenharia de Sistemas e Computação}
\linhapesquisa{Engenharia de Dados e Conhecimento}
\tipoproducao{bibliografica}
\projetovinculado{nao}
\agenciafomento{Conselho Nacional de Desenvolvimento Científico e
                Tecnológico}{CNPq}

\begin{document}
  \maketitle
  \frontmatter
  \begin{abstract}
    Resumo de teste, composto com serifa.
  \end{abstract}
  \begin{foreignabstract}
    Test abstract, set with serifs.
  \end{foreignabstract}
  \tableofcontents
  \mainmatter

  \chapter{Introdução}
    Primeiro parágrafo, para conferir o recuo e a família da fonte.%
    \footnote{A nota de rodapé sai em corpo menor, 2.2(b).}
  \section{Uma seção}
    Texto da seção.
    \begin{longquote}
      Citação de mais de três linhas, que sai em corpo menor e recuada,
      como manda a 2.2(b), e que aqui serve para conferir se o recuo e o
      corpo sobrevivem à troca de família da fonte padrão do documento.
    \end{longquote}
  \subsection{Uma Subseção}
    Texto terciário.
  \subsubsection{Uma Subsubseção}
    Texto quaternário.
  \paragraph{Um quinto nível}
    Texto quinário. Uma fórmula, $e^{i\pi} + 1 = 0$, para lembrar que a
    matemática não acompanha a família do texto.
\end{document}
%%
%%
%% End of file `tests/test_comserifa.tex'.
